% GENERATED FILE. Edit canonical lessons, then run npm run generate:books.
% canonical-source-hash: fnv1a64:bef1ffab5c23d767
% canonical-lessons: SA-C106-rajjuh, SA-C106-godhumah, SA-C106-tandulah, SA-C106-mudgah, SA-C106-supah

\chapter{Rope, Wheat, Rice, Mung beans, Dal}
\label{ch:sa-rope-wheat-rice-mung-beans-dal}
\hlchaptermodality{106}

\begin{chapteropening}
\textbf{By the end of this chapter:} \emph{I can say a rope, wheat, rice, mung beans and dal.}

Complete the last lesson of chapter 106: i can say a rope, wheat, rice, mung beans and dal.
\end{chapteropening}

\section[\texorpdfstring{\sk{रज्जुः} (rajjuḥ)}{रज्जुः (rajjuḥ)}]{\sk{रज्जुः} (rajjuḥ) — a rope}

\label{lesson:SA-C106-rajjuh}



\begin{quote}
Before the new one: say the Sanskrit for a fan, then the Sanskrit for a box.
\end{quote}

\subsection*{You'll want to know: \sk{रज्जुः}}
\textbf{\sk{रज्जुः}} — \emph{rajjuḥ} — \textquotedblleft{}a rope\textquotedblright{}.

\begin{grammarlens}[title={What you've built}]
One more word for this chapter.
\end{grammarlens}

\subsection*{Your turn}
\begin{itemize}
\raggedright
  \item \emph{Say it:} \emph{rajjuḥ}
  \item \emph{Say it:} \emph{rajjuḥ}, once more
  \item \emph{Say it:} say \emph{rajjuḥ}
  \item \emph{Recall:} say the Sanskrit for a fan, then the Sanskrit for a box, then say \emph{rajjuḥ} again
\end{itemize}

\begin{tcolorbox}[breakable,colback=teal!4,colframe=teal!35!black,title={Before you move on}]
What does \sk{रज्जुः} mean? (\textquotedblleft{}A rope\textquotedblright{}.) Say it once more.
\end{tcolorbox}
\section[\texorpdfstring{\sk{गोधूमः} (godhūmaḥ)}{गोधूमः (godhūmaḥ)}]{\sk{गोधूमः} (godhūmaḥ) — wheat}

\label{lesson:SA-C106-godhumah}



\begin{quote}
Before the new one: say the Sanskrit for a box, then the Sanskrit for a rope.
\end{quote}

\subsection*{You'll want to know: \sk{गोधूमः}}
\textbf{\sk{गोधूमः}} — \emph{godhūmaḥ} — \textquotedblleft{}wheat\textquotedblright{}.

\begin{grammarlens}[title={What you've built}]
One more word for this chapter.
\end{grammarlens}

\subsection*{Your turn}
\begin{itemize}
\raggedright
  \item \emph{Say it:} \emph{godhūmaḥ}
  \item \emph{Say it:} \emph{godhūmaḥ}, once more
  \item \emph{Say it:} say \emph{godhūmaḥ}
  \item \emph{Recall:} say the Sanskrit for a box, then the Sanskrit for a rope, then say \emph{godhūmaḥ} again
\end{itemize}

\begin{tcolorbox}[breakable,colback=teal!4,colframe=teal!35!black,title={Before you move on}]
What does \sk{गोधूमः} mean? (\textquotedblleft{}Wheat\textquotedblright{}.) Say it once more.
\end{tcolorbox}
\section[\texorpdfstring{\sk{तण्डुलः} (taṇḍulaḥ)}{तण्डुलः (taṇḍulaḥ)}]{\sk{तण्डुलः} (taṇḍulaḥ) — rice (uncooked)}

\label{lesson:SA-C106-tandulah}



\begin{quote}
Before the new one: say the Sanskrit for a rope, then the Sanskrit for wheat.
\end{quote}

\subsection*{You'll want to know: \sk{तण्डुलः}}
\textbf{\sk{तण्डुलः}} — \emph{taṇḍulaḥ} — \textquotedblleft{}rice (uncooked)\textquotedblright{}.

\begin{grammarlens}[title={What you've built}]
One more word for this chapter.
\end{grammarlens}

\subsection*{Your turn}
\begin{itemize}
\raggedright
  \item \emph{Say it:} \emph{taṇḍulaḥ}
  \item \emph{Say it:} \emph{taṇḍulaḥ}, once more
  \item \emph{Say it:} say \emph{taṇḍulaḥ}
  \item \emph{Recall:} say the Sanskrit for a rope, then the Sanskrit for wheat, then say \emph{taṇḍulaḥ} again
\end{itemize}

\begin{tcolorbox}[breakable,colback=teal!4,colframe=teal!35!black,title={Before you move on}]
What does \sk{तण्डुलः} mean? (\textquotedblleft{}Rice (uncooked)\textquotedblright{}.) Say it once more.
\end{tcolorbox}
\section[\texorpdfstring{\sk{मुद्गः} (mudgaḥ)}{मुद्गः (mudgaḥ)}]{\sk{मुद्गः} (mudgaḥ) — mung beans}

\label{lesson:SA-C106-mudgah}



\begin{quote}
Before the new one: say the Sanskrit for wheat, then the Sanskrit for rice.
\end{quote}

\subsection*{You'll want to know: \sk{मुद्गः}}
\textbf{\sk{मुद्गः}} — \emph{mudgaḥ} — \textquotedblleft{}mung beans\textquotedblright{}.

\begin{grammarlens}[title={What you've built}]
One more word for this chapter.
\end{grammarlens}

\subsection*{Your turn}
\begin{itemize}
\raggedright
  \item \emph{Say it:} \emph{mudgaḥ}
  \item \emph{Say it:} \emph{mudgaḥ}, once more
  \item \emph{Say it:} say \emph{mudgaḥ}
  \item \emph{Recall:} say the Sanskrit for wheat, then the Sanskrit for rice, then say \emph{mudgaḥ} again
\end{itemize}

\begin{tcolorbox}[breakable,colback=teal!4,colframe=teal!35!black,title={Before you move on}]
What does \sk{मुद्गः} mean? (\textquotedblleft{}Mung beans\textquotedblright{}.) Say it once more.
\end{tcolorbox}
\section[\texorpdfstring{\sk{सूपः} (sūpaḥ)}{सूपः (sūpaḥ)}]{\sk{सूपः} (sūpaḥ) — dal, soup}

\label{lesson:SA-C106-supah}



\begin{quote}
Before the new one: say the Sanskrit for rice, then the Sanskrit for mung beans.
\end{quote}

\subsection*{You'll want to know: \sk{सूपः}}
\textbf{\sk{सूपः}} — \emph{sūpaḥ} — \textquotedblleft{}dal, soup\textquotedblright{}.

\begin{grammarlens}[title={What you've built}]
One more word for this chapter.
\end{grammarlens}

\subsection*{Your turn}
\begin{itemize}
\raggedright
  \item \emph{Say it:} \emph{sūpaḥ}
  \item \emph{Say it:} \emph{sūpaḥ}, once more
  \item \emph{Say it:} say \emph{sūpaḥ}
  \item \emph{Recall:} say the Sanskrit for rice, then the Sanskrit for mung beans, then say \emph{sūpaḥ} again
\end{itemize}

\begin{tcolorbox}[breakable,colback=teal!4,colframe=teal!35!black,title={Before you move on}]
What does \sk{सूपः} mean? (\textquotedblleft{}Dal, soup\textquotedblright{}.) Say it once more.
\end{tcolorbox}
